\subsection{Revealing-optimization instance optimality}
\label{sec:common-upper-instance}

We first specialize the common transfer to revealing optimization.  Here a
direct first touch is raw execution: it always occupies the machine for the
public time $u$, independently of the hidden value.  Thus $r(p)=u$, and the
only model-specific work is to expand the fluid block area into an explicit
quadratic in the tested fraction.

\begin{theorem}[Revealing-optimization instance optimality]
\label{thm:common-upper-instance}
For every fixed $u>0$ there exist randomized nonanticipating algorithms
$\mathcal A^*_{n,u}$, depending only on $n,u$, and a deterministic sequence
$\eps_u(n)\to0$ such that every vector $p\in[0,u]^n$ satisfies
\begin{equation}
 \EE\ALG_{\mathcal A^*_{n,u}}(p)
 \le n^2\PhiRO(D[p])+\eps_u(n)n^2.
 \label{eq:cui-upper}
\end{equation}
Conversely, every randomized nonanticipating algorithm $\mathcal A$, even an
announced one, satisfies
\begin{equation}
 \EE\ALG_{\mathcal A}(p)
 \ge n^2\PhiRO(D[p])-\eps_u(n)n^2.
 \label{eq:cui-lower}
\end{equation}
One may take
\begin{equation}
 \eps_u(n)=O_u\!\left(n^{-1/6}\sqrt{\ln(n+2)}\right).
 \label{eq:cui-rate}
\end{equation}
\end{theorem}

\subsubsection{The explicit revealing benchmark}
\label{sec:cui-fluid}

Let $D$ be a probability distribution on $[0,u]$.  Apply
\cref{lem:maximum-density-module}, choose its threshold selector $g_*$, and
write
\begin{equation}
 a=a_D(g_*),\qquad
 \ell=\int_{(0,u]}p g_*(p)\,dD(p),
 \qquad \tau_D=\frac{1+\ell}{a}.
 \label{eq:cui-tau}
\end{equation}
Let $\nu$ be the residual measure from
\cref{eq:common-module-data}.  It has mass $1-a$, first moment
\begin{equation}
 \ell_\nu=\int p\,d\nu(p)=\int p\,dD(p)-\ell,
 \qquad
 \SPT(\nu)=\frac12\iint\min\{p,p'\}\,d\nu(p)d\nu(p'),
 \label{eq:cui-residual-moments}
\end{equation}
and support in $[\tau_D,u]$.

For a fixed tested fraction $q$, revealing optimization is the common fluid
model with $v=u$.  Hence \cref{lem:test-module-envelope} gives the blocks
\[
 qa\delta_{\tau_D}+q\nu+(1-q)\delta_u.
\]
When $\tau_D<u$, shortest-first order runs the testing block, the residual
types in SPT order, and finally the raw block.  Expanding
\cref{eq:common-envelope-area} gives
\begin{equation}
 F_{\mathsf{RO},u,D}(q)
 =(1+\ell)\left(q-\frac{aq^2}{2}\right)
 +q\ell_\nu(1-q)+q^2\SPT(\nu)
 +\frac{u(1-q)^2}{2}.
 \label{eq:cui-F}
\end{equation}
The four terms are, respectively, the testing block, the residual work
delaying raw jobs, the internal residual SPT cost, and the raw block.

If $\tau_D\ge u$, every testing or residual block has work per completion at
least $u$, so pure raw execution pointwise dominates every positive tested
fraction.  We therefore define
\begin{equation}
 \PhiRO(D)=
 \begin{cases}
  \displaystyle\min_{0\le q\le1}F_{\mathsf{RO},u,D}(q),&\tau_D<u,\\[3pt]
  u/2,&\tau_D\ge u.
 \end{cases}
 \label{eq:cui-Phi}
\end{equation}

The minimization is explicit.  Expanding gives
\begin{equation}
 F_{\mathsf{RO},u,D}(q)=\frac u2+Aq+Bq^2,
 \quad
 A=1+\int p\,dD(p)-u,
 \quad
 B=-\frac{a(1+\ell)}2-\ell_\nu+\SPT(\nu)+\frac u2.
 \label{eq:cui-quadratic}
\end{equation}
For each work budget, convexly combining feasible states for two tested
fractions gives a feasible state for their convex combination.  Thus the
completion envelope is concave in $q$, its unfinished-mass area is convex,
and $B\ge0$.  Moreover,
\[
 \tau_D<u
 \quad\Longleftrightarrow\quad
 \int(u-p)\,dD(p)>1
 \quad\Longleftrightarrow\quad A<0.
\]
Consequently a minimizing tested fraction is
\begin{equation}
 q^*=1\quad(B=0),
 \qquad
 q^*=\min\{1,-A/(2B)\}\quad(B>0).
 \label{eq:cui-qstar}
\end{equation}

\begin{proof}[Proof of \cref{thm:common-upper-instance}]
Take $r(p)=u$ and $Q=[0,1]$ in
\cref{thm:common-fluid-instance-transfer}.  Then
$v_r(D)=u$, and \cref{eq:cui-F,eq:cui-Phi} show that the common value
$\Phi_{r,Q}(D)$ is exactly $\PhiRO(D)$.  The upper and lower conclusions are
\cref{eq:cui-upper,eq:cui-lower}, and
\cref{eq:common-transfer-rate} gives \cref{eq:cui-rate}.
\end{proof}
